\documentclass[12pt]{article}
\usepackage[utf8]{inputenc}
\usepackage[T1]{fontenc}
\usepackage{hyperref}
\title{Persian Language Resources for K-12 Curriculum Development}
\author{Compiled by Claude Sonnet 5.5 (claude-sonnet-5-5) for Andr\'as Kornai}
\date{}
\begin{document}
\maketitle

\section{Introduction}
The code \texttt{fas} is the ISO 639-3 macrolanguage code for Persian (Farsi). The brief gave no name or identification, so this report assumes Persian. It is a large national language with abundant resources; the single search run here returned only a few credible items, so this is a short, incomplete list.

\section{Primary Educational Resources}
\begin{itemize}
\item Listing for \emph{Ketab Farsi Haftom}, described as Iran's seventh-grade Persian textbook; hosted on an unfamiliar site and not verified against Ministry of Education sources: \url{https://wiki.nektra.com/Download_PDFS/u35D7C/244880/Ketab Farsi Haftom.pdf}
\item \emph{Perfect Your Persian 3}, Chicago Persian School (heritage-learner series, elementary to advanced): \url{https://www.lulu.com/shop/chicago-persian-school/perfect-your-persian-3/paperback/product-nvmmmv7.html}
\end{itemize}

\section{Linguistic and Reference Materials}
\begin{itemize}
\item \emph{Modern Persian, Elementary Level}, Anthem Press (university textbook with companion site, audio and video): \url{https://anthempress.com/books/modern-persian-elementary-level-pb}
\end{itemize}

\section{Religious and Community Resources}
None found in this search.

\section{Digital and Audio Resources}
The Anthem Press title above is described as having a companion site with audio and video. No other digital resources were found.

\section{Curriculum Development Context}
Iranian national curriculum textbooks are the natural core for native-language K-12 use, but authoritative sources (Ministry of Education publishers) were not located in this search and should be checked directly. Afghan Dari and Tajik varieties are covered by other codes and were not searched.
\end{document}